\section{From Borrowed Answers to Earned Intelligence}
\label{sec:earned-intelligence-chain}

This development sequence tested a specific operational distinction between
\emph{answer access} and \emph{earned capability}.  Answer access denotes success
that depends on an open source containing a task-specific solution.  Earned
capability denotes a verified method that remains usable after its original
evidence is closed, transfers to a source-disjoint problem, survives restart,
and remains subordinate to independent outcome authority.  The resulting stack
progresses from a controlled learning differential to method invention,
verification-contract invention, local tool acquisition, and finally the
discovery and delayed re-observation of remote public authorities:
\[
\begin{aligned}
&\text{retained method}\rightarrow
\text{invented explanation}\rightarrow
\text{invented verifier}\rightarrow
\text{acquired adapter}\\
&\qquad\rightarrow
\text{shortlist-free authority discovery}\rightarrow
\text{precommitment}\rightarrow
\text{later consequence}\rightarrow
\text{retention}.
\end{aligned}
\]
Every promoted procedure preserved the separation between proposal and
authority: a plausible explanation, tool, endpoint, or response format did not
become knowledge until it survived the relevant counterexample, execution,
transfer, safety, and delayed-outcome gates.

\subsection{Controlled earned-intelligence differential}

The first experiment compared three systems over six method families:
integrity, provenance, mathematics, transactions, causal investigation, and
planning.  The \emph{retained} system could use methods learned from verified
experience; the \emph{cold} system received the same candidate-method space but
no retained selection history; and the \emph{answer-book} control received the
correct task-specific method only while its source remained open.  Retained
AION succeeded immediately on all six families and again after source closure.
The cold system eventually solved all six but required three attempts per
family.  The answer-book system scored $6/6$ while the source was available and
$0/6$ after closure.  The corresponding reduction in investigation attempts
was
\[
\Delta_{\mathrm{attempt}}
=1-\frac{A_{\mathrm{warm}}}{A_{\mathrm{cold}}}
=66.67\%.
\]
Thus, all six families exhibited a positive retained-learning differential.
The initial live result was $5/6$ because a causal method card referenced an
unavailable artifact name.  The run was rejected; the gate was not weakened.
The procedure located the actual promoted causal artifact, rebound the evidence,
and passed the unchanged cohort.  HexCore then promoted
\texttt{procedure\_earned\_intelligence\_differential\_arena\_v1}.

\begin{table}[t]
\centering
\caption{Controlled distinction between retained competence and temporary
answer access.}
\label{tab:earned-intelligence-differential}
\begin{tabular}{lccc}
\toprule
System & Source open & Source closed & Attempts per family \\
\midrule
Retained AION & $6/6$ & $6/6$ & $1$ \\
Cold control & $6/6$ eventually & $6/6$ eventually & $3$ \\
Answer-book control & $6/6$ & $0/6$ & not applicable \\
\bottomrule
\end{tabular}
\end{table}

\subsection{Open explanation invention and falsification}

The next arena removed supplied relation and operator labels.  AION observed six
opaque three-input Boolean worlds and searched an executable expression
meta-grammar for candidate explanations.  For a version space
$\mathcal{H}_{t}$ consistent with observations $D_t$, the active experiment was
selected at a point of maximal predictive disagreement,
\[
x_t^{\star}
=\arg\max_x
\left|\left\{h(x):h\in\mathcal{H}_{t}\right\}\right|.
\]
The proposed prediction or abstention was committed before the withheld outcome
was exposed.  Four worlds supported immediate correct predictions.  One world
left two equally defensible explanations, causing a correct abstention rather
than a forced answer.  In a sixth world, the initially preferred simpler theory
made a confident but incorrect prediction; the independent counterexample
falsified it, after which the surviving theory was retained.  Across 24 active
queries, the revised method transferred to renamed schemas in all six worlds.
An unsupported four-variable problem produced an out-of-distribution
abstention.

The first challenger had treated a minimum-description-length winner as
epistemic certainty and achieved only $5/6$.  It was rejected.  The promoted
successor explicitly separated unique support, ambiguity, and falsification,
yielding zero forced predictions under ambiguity and zero unsafe actions.
HexCore promoted
\texttt{procedure\_open\_earned\_intelligence\_invention\_v1}.

\subsection{Heterogeneous verification-contract invention}

AION was then given five genuine previously completed North-Star projects, but
only their objectives, success statements, and artifact bytes.  Family labels,
original verifier metadata, hidden fields, and later outcome rows were withheld.
For each project AION constructed a typed verification contract
\[
K_i=(P_i,V_i,A_i,C_i),
\]
where $P_i$ is an observable property, $V_i$ an executable verifier, $A_i$ the
permitted outcome authority, and $C_i$ the counterexample that the contract must
reject.  The contract hash was committed before the corresponding historical
consequence was opened.

The five independently invented authority programs used different epistemic
semantics: integer recomputation for a mathematical witness; immutable digest
and exact-span recovery for a document claim; fresh-subprocess positive and
negative discrimination for executable behaviour; later public revision
comparison for repository research; and delayed numerical-envelope checking for
an environmental decision.  All five contracts accepted their independent
positive outcomes and rejected their adversarial mutations.  An unsupported
artifact format caused abstention.  There were no unsafe actions or live writes,
and the champion survived restart.  HexCore promoted
\texttt{procedure\_heterogeneous\_project\_contract\_invention\_v1}.

\subsection{Acquiring missing local execution adapters}

The verifier layer next confronted valid contracts for which no adapter was
registered.  In the first acquisition stage, AION received an engineered tool
search surface but no executable paths or argument templates.  It inspected
published local interfaces, synthesized typed read-only argument vectors,
security-scanned them, executed a development property, and transferred each
adapter to a source-disjoint artifact.  It acquired \texttt{xmllint} for
non-networked XML/SVG structural checks, \texttt{file} for media-type
classification, and \texttt{openssl dgst} for independent SHA-256 verification.
All three passed development and transfer tests; six malicious variants
requesting output files, signing keys, shell execution, network access,
unapproved executables, or write-enabled operation were rejected before
execution.  The retained policy reduced search attempts by $66.67\%$ and
required no relearning after restart.  This produced
\texttt{procedure\_open\_cli\_adapter\_acquisition\_v1}.

The successor removed both the tool shortlist and supplied requirement labels.
From three natural-language objectives, AION first formulated an observable
property and then searched 267 candidates in the operating system's public
manual index.  Candidate interfaces were ranked by semantic fit, provenance,
capability risk, and execution cost.  The selected adapters were \texttt{jq -e
.} for structural JSON parsing, \texttt{wc -l} for physical-line counting, and
\texttt{grep -F -c} for exact-literal matching.  Each result was checked by an
independent Python recomputation and transferred to a different artifact.

The first property-first challenger was rejected for two distinct errors: its
macOS help parser failed to understand compact flag groups and selected
\texttt{msggrep} rather than \texttt{grep}; and its trust policy considered
\texttt{/bin/sh} safe merely because the executable occupied a trusted path.
The revised challenger learned that provenance and capability class are
separate dimensions, added grouped-flag parsing and capability-level denial,
and passed without weakening a safety gate.  Results were $3/3$ property
formulations, $3/3$ development checks, $3/3$ source-disjoint transfers, $6/6$
malicious-candidate rejections, correct ambiguous-objective abstention, and zero
network actions or live writes.  HexCore promoted
\texttt{procedure\_property\_first\_adapter\_discovery\_v1}.

\subsection{Shortlist-free remote authority acquisition}

The next stage replaced local interfaces with an unfamiliar remote scientific
mission.  AION received a broad objective concerning recent earthquake risk,
with no source catalogue, documentation address, or endpoint.  It generated a
search query, inspected ten public results, rejected commercial and
non-government candidates, recovered two official U.S. Geological Survey
documents, and extracted 21 callable endpoints.  It selected the official USGS
GeoJSON feeds and constructed a credential-free HTTPS \texttt{GET} adapter whose
observables were event count, maximum magnitude, and latest event timestamp.
The required response contract contained \texttt{bbox}, \texttt{features},
\texttt{metadata}, and \texttt{type}.  Development on the hourly feed
transferred to the daily feed, while five illegitimate remote candidates were
rejected.

This success followed a meaningful rejection chain.  The configured OpenAI
teacher returned HTTP~429, and two local-model proposals hallucinated
unsupported government API roots.  None was treated as authority.  Grounded
search plus official-document recovery supplied the successful path.  During
the promotion run, the live feeds contained six hourly events and 221 daily
events; these values were retained as time-indexed evidence rather than timeless
claims.

Before the later request, AION durably committed the schema and physical bounds.
After a five-minute wall-clock boundary, the continuous outcome learner issued a
fresh request.  The response digest had changed, while the committed schema and
bounds remained valid.  HexCore promoted both
\texttt{procedure\_shortlist\_free\_remote\_authority\_acquisition\_v1} and its
delayed child,
\texttt{procedure\_later\_confirmed\_remote\_authority\_retention\_v1}.

\subsection{Cross-domain transfer of authority acquisition}

The grounded discovery method was finally transferred to two unrelated
missions, again without endpoint catalogues.  For public finance, AION recovered
U.S. Treasury Fiscal Data documentation and a read-only response contract over
the \texttt{data}, \texttt{links}, and \texttt{meta} containers.  For software
security, the first live challenger selected NIST's \texttt{manifest.json}.
Although the file was official, safely accessible, and valid JSON, it contained
no vulnerability evidence.  The run was rejected, establishing
\[
\text{official provenance}+\text{valid syntax}
\not\Rightarrow \text{mission-relevant authority}.
\]
The revised challenger required semantic evidence for the target property and
selected CISA's Known Exploited Vulnerabilities catalogue, which contained
1,656 records at the time of acquisition.

The original USGS mission had required four proposal paths: a rate-limited
teacher, two unsupported local proposals, and one grounded-search path.  Both
new domains selected grounded discovery on the first path.  Relative to replaying
the parent cost in both domains, the acquisition-attempt reduction was
\[
1-\frac{1+1}{4+4}=75\%.
\]
Both contracts were precommitted and independently reissued after five-minute
boundaries.  Their response hashes happened to remain unchanged during that
short interval, but both fresh observations preserved the committed containers
and nonnegative record properties.  Consequently, this result demonstrates
independently timed re-observation and retained contract validity, not a claim
that the external worlds changed.  HexCore promoted
\texttt{procedure\_cross\_domain\_remote\_authority\_transfer\_v1} and
\texttt{procedure\_cross\_domain\_remote\_later\_retention\_v1}.

\begin{table*}[t]
\centering
\caption{Integrated results for the earned-intelligence development chain.}
\label{tab:earned-intelligence-chain}
\begin{tabular}{p{0.27\textwidth}p{0.12\textwidth}p{0.14\textwidth}p{0.33\textwidth}}
\toprule
Capability & Positive result & Adversarial result & Principal evidence \\
\midrule
Retained-method differential & $6/6$ & answer book $0/6$ closed & $66.67\%$ fewer attempts than cold \\
Open explanation invention & $6/6$ transfer & ambiguity and OOD abstention & 24 active queries; one falsified theory revised \\
Verifier-contract invention & $5/5$ & $5/5$ mutations rejected & five distinct authority programs \\
Open CLI acquisition & $3/3$ transfer & $6/6$ malicious rejected & $66.67\%$ attempt reduction \\
Property-first discovery & $3/3$ transfer & $6/6$ malicious rejected & 267 interfaces considered; no shortlist \\
Remote scientific authority & hourly-to-daily pass & $5/5$ remote candidates rejected & later changed USGS digest retained \\
Cross-domain remote transfer & $2/2$ & irrelevant official JSON rejected & Treasury and CISA; $75\%$ attempt reduction \\
\bottomrule
\end{tabular}
\end{table*}

\subsection{Persistent operation, integrity, and claim boundary}

The delayed challenges are handled by the continuous real-outcome learner,
installed as a macOS-supervised service with automatic restart.  Pending
commitments remain durable, eligible challenges are closed only after their
wall-clock boundaries, and observations retain their response hashes,
authorities, and verdicts.  The complete focused earned-intelligence stack
passed $7/7$ regression tests.  Across the chain there were no credential uses,
non-\texttt{GET} remote actions, unsafe executions, or live-repository writes.

These results support a bounded claim: AION can distinguish retained method
learning from temporary answer access; invent and revise explanations within an
engineered meta-grammar; infer heterogeneous verification semantics from prior
real project artifacts; acquire safe local adapters; discover official remote
authorities without endpoint catalogues; and retain those authorities only
after independently timed consequences.  They do not establish unrestricted
web agency, arbitrary verifier or tool invention, complete domain mastery,
Autonomous General Apprentice status, or AGI.  Boolean and property grammars,
cohort selection, trust rules, authority adapters, query interpretation, and
promotion gates remain partly engineered.  The next removal target is
multi-step document and institutional intelligence in which AION must discover
intermediate identifiers, follow dependent information actions, and construct
a verifier when no documentation page exposes an immediately executable
endpoint.  Advanced/Expert competency, natural self-repair, and week-scale
retention remain separate evidence gates and receive no credit from this chain
alone.
